\chapter{Successioni e funzioni continue}
	\pagestyle{plain}
	\thispagestyle{empty}
	\pagestyle{fancy}
In questo capitolo andremo a definire il concetto di successione e vedere il forte collegamento presente fra esse e il concetto di continuità in più variabili. \\
Ricordiamo al lettore che una successione a valori in un insieme $A$ è una funzione $x_k: \mathbb{N} \to A$. Vediamo adesso come si definisce il concetto di \emph{successione convergente} in più variabili

%
%
%
\vspace{0.5cm}
\section{Convergenza e limite di una successione}

\begin{definition}[\textcolor{blue}{convergenza di una successione}]
diremo che $\{x_k \} \subset \mathbb{R}^n$ è una successione convergente a $z \in \mathbb{R}^n$ se $\lim\limits_{k \to +\infty} |x_k - z| = 0$.
\end{definition}

\begin{remark}
Naturalmente se $\underline{x_k} = (x_{k_1}, x_{k_2}, \ldots x_{k_n}) \in \mathbb{R}^n$ converge a $\underline{z} = (z_1, z_2, \ldots z_n)$ 
\[\iff \lim\limits_{k \to \infty} x_{k_i} = z_i \quad \forall\ i \in 1, \ldots n\]
\end{remark}

\begin{example}
Consideriamo la successione $\underline{x_k} = (e^{-k} + 1,\frac{(-1)^k}{k})$. Per il precedente teorema questa successione è convergente.
\end{example}

%
\vspace{0.5cm}
\begin{prop}[\textcolor{blue}{Unicità del limite}]
Il limite di una successione è unico. 
Se $\underline{x_k} \in \mathbb{R}^n \, \, \forall k$ converge a $\underline{z} \in \mathbb{R}^n$ allora $z$ è unico.
\end{prop}

\begin{proof}
Supponiamo per assurdo che $\underline{x_k}$ converga a $\underline{z}$ e $\underline{y}$ con $\underline{z} \neq \underline{y}$. Allora
$$
|\underline{z} - \underline{y}|\ \leq\ |\underline{x_k} - \underline{z}|\, +\, |\underline{x_k} - \underline{y}|\ \stackrel{k \to +\infty}{\to}\ 0\ \implies \ z\, =\, y\ \ .
$$
\end{proof}

%
\vspace{0.5cm}
\begin{prop}
Se $\underline{x_k} \to \underline{x}$ allora $|\underline{x_k}| \to |\underline{x}|$
\end{prop}

\begin{proof}
si osserva che, dalla proposizione~\ref{prop:dis_triang}, si ha che
$$
|\underline{x_k}| - |\underline{x}| \leq |\underline{x_k} - \underline{x}| \stackrel{k \to +\infty}{\to} 0
$$
\end{proof}

\vspace{0.5cm}
\noindent Mostriamo adesso alcune proposizioni che ci permettono di comprendere alcune proprietà con cui trattare queste successioni convergenti, come ad esempio come fare operazioni fra successioni e cosa questo comporta come risultato per quanto riguarda i punti a cui convergono.

\begin{prop}[\textcolor{blue}{spazio vettoriale delle successioni convergenti}]
Se $\underline{x_k} \to \underline{x}$ e $\underline{y_k} \to \underline{y}$ allora $\forall \lambda, \mu \in \mathbb{R}, \, \lambda \underline{x_k} + \mu \underline{y_k} \to \lambda \underline{x} + \mu \underline{y}$ 
\end{prop}

\begin{proof}
	osserviamo che
	$$
	|\lambda \underline{x_k} + \mu \underline{y_k} - \lambda \underline{x} - \mu \underline{y}| = |\lambda (\underline{x_k} - \underline{x}) + \mu (\underline{y_k} - \underline{y})|
	$$
	Per la proposizione~\ref{prop:dis_triang} abbiamo che
	$$
	|\lambda (\underline{x_k} - \underline{x}) + \mu (\underline{y_k} - \underline{y})| \leq |\lambda| |\underline{x_k} - \underline{x}| + |\mu| |\underline{y_k} - \underline{y}| \stackrel{k \to +\infty}{\to} 0
	$$
	dunque la tesi
\end{proof}

\begin{remark}
l'importanza di questa dimostrazione sta nel fatto che questo teorema dimostra che l'insieme delle successioni convergenti in $\mathbb{R}^n$ forma uno spazio vettoriale che è chiuso rispetto all'addizione $+_{\mathbb{R}}$ e prendendo come prodotto per scalare $*_{\mathbb{R}}$
\end{remark}


\begin{prop}
	Dato un vettore $\underline a\in \mathbb{R}^n$ e una successione $\underline{x_k}\rightarrow \underline{x_\infty}$ in $\mathbb{R}^n$, allora $\rightarrow \underline{x_k}\cdot \underline a \rightarrow \underline{x_\infty}\cdot \underline a$ .
\end{prop}

\begin{proof}
	\[ \lim_{k\rightarrow\infty}\, |\underline x_k\cdot\underline a - \underline x_\infty\cdot\underline a|\ =\ \lim_{k\rightarrow\infty}\, |(\underline x_k - \underline x_\infty) \cdot\underline a|\ \leq\ \lim_{k\rightarrow\infty}\, |\underline x_k - \underline x_\infty|\cdot |\underline a|\ =\ 0 \]
\end{proof}


%
%
%
\vspace{0.5cm}
\section{Successioni e topologia}

\noindent Adesso andiamo a mostrare una proprietà che segue direttamente dalla topologia di $\mathbb{R}^n$ (in spazi topologici qualunque non è sempre valido)

\begin{prop}[\textcolor{blue}{caratterizzazione degli insiemi chiusi}]
	Un insieme $A \subset \mathbb{R}^n$ è chiuso se e solo se $\forall \underline{x_k} \to \underline{z}, \underline{x_k} \in A \, \forall k$ allora $z \in A$. Cioè se ogni successione contenuta in $\mathbb{R}^n$ converge ad un punto $\underline z$ di $\mathbb{R}^n$. 
	
	Formalmente:
	$$
	(A \subset \mathbb{R}^n \text{ è chiuso}) \iff (\forall\ \underline{x_k} \to \underline{z},\ \underline{x_k} \in A \ \forall k \implies z \in A) 
	$$
	\label{prop:caratt_chiusi}
\end{prop}

\begin{proof} \hspace{1em} \newline
$\boxed{\Rightarrow}$: Se $A$ è chiuso, considerando una generica $\underline{x_k} \to \underline{z}, \underline{x_k} \in A \, \forall k$, allora $\forall r > 0$ dato che $|\underline{x_k} - \underline{z}| \to 0 \, \exists k_r \in \mathbb{N}$ tale che $|\underline{x_k} - \underline{z}| < r \, \forall k \geq k_r \implies \underline{x_k} \in B(z, r) \implies B(z, r) \cap A \neq \emptyset \implies z \in \bar{A}$. Siccome A è chiuso allora $\bar{A} = A$ dunque $z \in A$. \\
$\boxed{\Leftarrow}$: Sia $\underline{w} \in \bar{A} \implies B(\underline{w}, \frac{1}{k}) \cap A \neq \emptyset \, \forall k \geq 1$ quindi esiste una successione $\underline{x_k}$ (potremmo prendere per esempio $\underline{w}-\frac{1}{k}$) tale che $|\underline{x_k} - \underline{w}| < \frac{1}{k} \implies \underline{x_k} \stackrel{k \to +\infty}{\to} \underline{w} \implies \underline{w} \in A$ per la seconda proprietà, ma allora segue che $\bar{A} = A$, dunque $A$ è chiuso per definizione. 
\end{proof}

\noindent Questa proprietà è molto importante, siccome garantisce l'equivalenza fra \textbf{chiusura sequenziale} e \textbf{chiusura} (che negli spazi topologici non è generalmente garantito) in $\mathbb{R}^n$. \\
In $\mathbb{R}^n$, dato un insieme $A \subset \mathbb{R}^n$, se vogliamo verificare che sia chiuso sarà dunque necessario verificare che ogni successione $\underline{x_k}$ tali che $\forall k, \underline{x_k} \in A$ avremo che $\underline{x_k} \to \underline{x} \in A$, ovvero che ogni successione convergente in $A$ converga ad un punto che appartiene ad $A$.

%
%
%
\vspace{0.5cm}
\section{Funzioni continue}

\begin{definition}[\textcolor{blue}{funzioni continue}]
Una funzione $f: A \to \mathbb{R}$ con $A \subset \mathbb{R}^m$ si dice continua in $\underline{z} \in A$ se $$\forall \underline{x_k} \in A, \underline{x_k} \to \underline{z}, \, \lim\limits_{k \to +\infty} f(\underline{x_k}) = f(\underline{z})$$
\end{definition}

\begin{remark}
Diremo che $f: A \to \mathbb{R}$ con $A \subset \mathbb{R}^n$ è continua in $A$ se è continua $\forall \underline{x} \in A$.
\end{remark}

\noindent Andiamo a ricordare un concetto molto importante preso dal precedente capitolo:

\vspace{0.5cm}
\begin{definition}[\textcolor{blue}{apertura relativa}]
Fissato $A \subset \mathbb{R}^n$ diremo che $S \subset A$ è aperto in $A$ se $\exists \Omega \subset \mathbb{R}^n$ aperto tale che $$S = A \cap \Omega$$
\end{definition}
Dopo questa definizione siamo pronti per enunciare una serie di teoremi sulle funzioni continue che le caratterizzano in $\mathbb{R}^n$.

\vspace{0.5cm}
\begin{theorem}[C1]
Sia $f: A \to \mathbb{R}^m$ con $A \subset \mathbb{R}^n$. Allora $f$ è continua $\iff \forall U \subset \mathbb{R}^m$ aperto $f^{-1}(U)$ è aperto in $A$
\label{thm:theo_c1}
\end{theorem}

\noindent \textbf{Dimostrazione. } 
$\boxed{\Rightarrow}$: se $f$ è continua, supponiamo per assurdo che esista $U \subset \mathbb{R}^m$ aperto tale che $f^{-1}(U)$ non è aperto in $A$. Allora $\exists z \in f^{-1}(U)$ tale che $\forall r > 0, B(\underline{z}, r) \cap A \not\subseteq f^{-1}(U)$, dunque $\forall k \geq 1, B(\underline{z}, \frac{1}{k}) \cap A \not\subseteq f^{-1}(U)$. Ma ciò implica che $\exists \underline{x_k} \in A \cap B(\underline{z}, \frac{1}{k})$ tale che $\underline{x_k} \not\in f^{-1}(U)$. Ma allora 
$$
|\underline{x_k} - \underline{w}| < \frac{1}{k} \to 0 \implies f(\underline{x_k}) \to f(\underline{z})
$$
e chiaramente abbiamo che $f(x_k) \in U^c$ (se appartenesse a $U$ allora $\underline{x_k} \in f^{-1}(U)$ il che contraddice come abbiamo definito la successione). \\ Ma allora $f(z) \not\in U$: infatti, $U^c$ è chiuso dunque, per la proposizione~\ref{prop:caratt_chiusi}, abbiamo che ogni successione a valori in $U^c$ convergente ad un punto implica che tale punto appartenga all'insieme chiuso, in questo caso $U^c$, ma ciò contraddice l'ipotesi iniziale che $z \in f^{-1}(U)$. \\
$\boxed{\Leftarrow}$: sappiamo che la preimmagine di aperti è aperta in $A$ allora sia $\underline{x_k} \to \underline{z}$ con $\underline{x_k} \in A$ e $\underline{z} \in A$: se supponiamo, per assurdo, $f(\underline{x_k}) \not\to f(\underline{z})$ allora $\exists \alpha: \mathbb{N} \to \mathbb{N}$ tale che $\alpha(k) \stackrel{k \to +\infty}{\to} +\infty$ ed esiste $\sigma > 0: |f(\underline{x}_{\alpha(k)}) - f(\underline{z})| \geq \sigma$. D'altra parte $f^{-1}(B(f(z), \sigma))$ è aperto in $A$ (per ipotesi, essendo preimmagine di un aperto), dunque esiste $\delta > 0$ tale che $A \cap B(z, \delta) \subset f^{-1}(B(f(z)), \sigma)$ ed esiste $k_0 \geq 1$ tale che $x_k \in A \cap B(z, \delta) \subseteq f^{-1}(B(f(z)), \sigma) \, \forall k \geq k_0 \implies |f(\underline{x_k}) - f(\underline{z})| < \sigma$ e questo contraddice la precedente disequazione $|f(\underline{x}_{\alpha(k)})-f(\underline{z})| \geq \sigma$ per $k$ sufficientemente grande.
\hfill $\square$

\vspace{0.5cm}
\begin{theorem}[\textcolor{blue}{della permanenza del segno}]
Sia $f: A \to \mathbb{R}$ continua in $z \in A$. Se $f(z) > 0$ allora $\exists \delta > 0$ tale che $\forall x \in A \cap B(z, \delta), f(x) > 0$
\end{theorem}

\noindent \textbf{Dimostrazione. }
Se neghiamo la tesi allora otteniamo che $\forall \delta > 0, \exists x \in A \cap B(z, \delta): f(x) < 0$. Ma allora deduciamo che $\exists \underline{x_k} \in A \cap B(z, \frac{1}{k}): f(x_k) <0$, ma $\underline{x_k} \to z$ e $f(x_k) \to f(z) \leq 0$ che rappresenta una contraddizione.
\hfill $\square$

\vspace{0.5cm}
\begin{theorem}[\textcolor{blue}{C2}]
Siano $A \subset \mathbb{R}^n$, $f, g: A \to \mathbb{R}^m$ e $h: B \to \mathbb{R}^k$ con $f(A) \subset B \subset \mathbb{R}^m$. Se $f$ e $g$ sono continue in $z \in A$ e $h$ è continua in $f(z) \in B$ allora valgono le seguenti proprietà:
\begin{enumerate}[label=\protect\circled{\arabic*}]
	\item $\forall \lambda, \mu \in \mathbb{R}$ abbiamo che $\lambda f + \mu g$ è continua in $z \in A$;
	\item Se $m = 1$, allora $fg$ è continua in $z$;
	\item Se $m=1$ e $g(z) \neq 0$ allora $\exists \delta > 0$ per cui $f/g$ è ben definita su $A \cap B(z, \delta)$ ed è continua in $z$;
	 \item la composizione $h \circ f$ è continua in $z$
\end{enumerate}
\end{theorem}

\noindent \textbf{Dimostrazione. }
Per mostrare la $\circled{1}$ osserviamo che, prendendo $x_k \to z$, abbiamo che
$$
\lambda f(x_k) + \mu g(x_k) \to \lambda f(z) + \mu g(z)
$$
Per mostrare la $\circled{2}$ osserviamo che
$$
x_k \to z \implies f(x_k)g(x_k) \to f(z)g(z)
$$
Per mostrare la $\circled{3}$ si osserva che possiamo supporre che $g(z) > 0$ senza perdita di generalità da cui segue che
$$
\exists \delta > 0: \forall x \in A \cap B(z, \delta), g(x) > 0 \implies \frac{f}{g}: B(z, \delta) \to \mathbb{R}
$$
che ci permette di concludere che $\frac{f}{g}$ sia ben definita e $\lim\limits_{k \to +\infty} \frac{f(x_k)}{g(x_k)} = \frac{f(z)}{g(z)}$ con $x_k \to z$. \\
La $\circled{4}$ segue banalmente dal fatto che presa $x_k \in A, x_k \to z$ allora
$$
\lim_{k \to +\infty} h(f(x_k)) = h(\lim_{k \to +\infty} f(x_k)) = h(f(z))
$$
per la continuità di $h$ in $f(z) \in f(A)$ e di $f$ in $z \in A$.
\hfill $\square$

\vspace{0.5cm}
\begin{prop}
Sia $f: A \to \mathbb{R}^m$, dove $f(x) = (f_1(x), f_2(x), \ldots f_m(x))$ e $f_j(x): A \to \mathbb{R} \, \forall j \in \{1, \ldots m\}$ con $A \subset \mathbb{R}^n$. Allora abbiamo che
	$$f \text{ è continua in } p \in A \iff f_j \text{ è continua in p} \in A \, \forall j \in \{1, \ldots m\}$$
\end{prop}

\noindent \textbf{Dimostrazione. }
Basta osservare che $f(x_k) \to f(p) \iff f_j(x_k) \to f_j(p) \, \forall j \in \{1, \ldots m \}$
\hfill $\square$

\begin{remark}
$f$ è continua in $A \iff f_j$ è continua in $A \, \forall j \in \{1, \ldots m \}$.
\end{remark}

\vspace{0.5cm}
\begin{example}
La funzione $|\cdot|: \mathbb{R}^n \to [0; +\infty)$ è continua siccome abbiamo mostrato che se $x_k \stackrel{k \to +\infty}{\to} x \implies |x_k| \stackrel{k \to +\infty}{\to} |x|$
\end{example}

\vspace{0.5cm}
\begin{example}
	La funzione $f(x, y, z) = x^2 - 3y + \sqrt{|xz|}, f: \mathbb{R}^3 \to \mathbb{R}$ è continua in quanto somma di funzioni continue: $\sqrt{|xz|}$ è continua in quanto composizione della funzione $\sqrt{t}, t \geq 0$ con $|\cdot|$ e $xz$ che sono tutte continue. 
\end{example}

\vspace{0.5cm}
\begin{example}
	$f: \mathbb{R}^2 \setminus \{ 0 \} \to \mathbb{R}$ con 
	$$
	f(x, y) = \frac{xy^3}{x^2 + y^4}
	$$
	è continua in quanto quoziente di funzioni continue.
\end{example}

\vspace{0.5cm}
Vogliamo adesso individuare delle \emph{buone} proprietà che le funzioni continue possiedono. Per fare ciò dobbiamo definire qualche altro concetto topologico:
\vspace{0.5cm}
\begin{definition}
	Un insieme $E \subseteq \mathbb{R}^n$ è limitato se $\exists R > 0 : E \subseteq B(0, R)$
\end{definition}

\begin{remark}
	Un insieme è illimitato $\iff$ $\exists \{x_k\} \in E \, \, \forall k : |x_k| \to +\infty$
\end{remark}

\begin{remark}
	Ogni palla aperta o chiusa è limitata
\end{remark}

\begin{remark}
	Se $B$ è limitato e $A \subseteq B \implies A$ è limitato.
\end{remark}

\vspace{0.5cm}
\begin{example}
$$
A = \{(x, y, z) \in \mathbb{R}^3: x^2 + y^2 + 1 < z < 10 \}
$$
è limitato. Infatti preso $(x, y, z) \in A \implies |(x, y, z)|^2 = x^2 + y^2 + z^2 = x^2 + y^2 + z^2 + 1 - 1 < z + z^2 - 1 < 10 + 100 - 1 = 109 \implies A \subseteq B(0, \sqrt{109})$ dunque è limitato. \\
Da questo esempio intuiamo che per verificare la limitatezza di un insieme non è necessario disegnarlo (e spesso è anche molto difficile).
\end{example}

\vspace{0.5cm}
\begin{example}
$$
A = \{(x, y) \in \mathbb{R}^2: x^2 y^2 < 1 \} = \{ (x, y) \in \mathbb{R}^2 : |xy| < 1 \}
$$
osserviamo che è illimitato siccome possiamo prendere $x_k = (\frac{1}{2k}, k)$ e osservare che $|\frac{1}{2k} \cdot k| = \frac{1}{2} < 1$ dunque $\forall k, x_k = (\frac{1}{2k}, k) \in A$ ma per $k \to +\infty, |x_k| = \sqrt{\frac{1}{4k^2} + k^2} \geq |k| \stackrel{k \to +\infty}{\to} +\infty$. 
Questo insieme, nel piano cartesiano, corrisponde ai punti che stanno sotto alla figura $|xy| = 1$

\begin{figure}[H]
	\centering
	\begin{tikzpicture}
		% Assi cartesiani
		\draw[->] (-3.1, 0) -- (3.1, 0) node[right] {$x$};
		\draw[->] (0, -3.1) -- (0, 3.1) node[above] {$y$};
	  
		% Iperbole nel primo e terzo quadrante: y = 1/x
		\draw[thick, black, domain=0.333:3, samples=100] plot (\x, {1/\x});
		\draw[thick, black, domain=0.333:3, samples=100] plot (-\x, {-1/\x});
	  
		% Iperbole nel secondo e quarto quadrante: y = -1/x
		\draw[thick, black, domain=0.333:3, samples=100] plot (-\x, {1/\x});
		\draw[thick, black, domain=0.333:3, samples=100] plot (\x, {-1/\x});
	  
		% Etichetta per |xy| = 1
		\node at (1.2, 0.2) {\small $|xy| = 1$};
	  
	  \end{tikzpicture}
	  \caption{Immagine dell'insieme A}
	  \label{fig:set_A}
\end{figure}
\end{example}

%
%
%
\vspace{0.5cm}
\section{Insiemi compatti e connessi per archi}
\begin{definition}[\textcolor{blue}{insiemi compatti}]
	Un insieme $E \subseteq \mathbb{R}^n$ è compatto se $\forall \{x_k \} \subset E, \exists x \in E \exists \alpha: \mathbb{N} \to \mathbb{N}$ strettamente crescente tale che
	$$
	x_{\alpha(k)} \to x \in E
	$$
\end{definition}

\vspace{0.5cm}
\begin{theorem}[\textcolor{blue}{di Heine-Borel}]
	Dato $E \subseteq \mathbb{R}^n$, allora 
	
	\begin{align*}
		E \subseteq \mathbb{R}^n \text{ è compatto } \iff E \text{ è chiuso e limitato.}
	\end{align*}
	
	\label{thm:heine_borel}
\end{theorem}

\noindent \textbf{Dimostrazione. }
$\boxed{\Rightarrow}$: sia $x_k \to x \implies$ (per l'ipotesi di compattezza) $\exists y \in E : \exists \alpha: \mathbb{N} \to \mathbb{N} : x_{\alpha(k)} \to y \in E \implies x=y \in E$ per unicità del limite. Se $E$ non fosse limitato $\implies \exists \{x_k \} \subset E : |x_k| \to +\infty$ e per la compattezza $\exists \beta: \mathbb{N} \to \mathbb{N}$ strettamente crescente tale che $x_{\beta(k)} \to x \in E$ ma per unicità del limite $|x_\beta(k)| \to |x_k| \to +\infty$ e questo contraddice il limite $+\infty$. \\
$\boxed{\Leftarrow}$: Se $\{ x_ k\} \subset E$ e $x_k = (x_{k_1}, \ldots, x_{k_n}) \implies x_{k_j}$ sono successioni reali limitate tali che $\exists x_{{\alpha_1(k)}_1} \to x_1, \exists x_{{\alpha_2(k)}_2} \to x_2, \ldots \implies 
\exists x_{\alpha_1 \ldots \alpha_n(k)} \to x = (x_1, \ldots, x_n)$ (il fatto di poter estrarre una successione convergente è conseguenza del teorema di Bolzano-Weierstrass) e $x \in E$ per la chiusura di $E$. 
\hfill $\square$

\vspace{0.5cm}
\noindent A questo punto ricordiamo la definizione di massimo e minimo locale o globale:
\begin{definition}[\textcolor{blue}{massimo globale}]
Sia $A \subseteq \mathbb{R}^n$ e sia $f: A \to \mathbb{R}$. Diremo che $z \in A$ è un punto di massimo assoluto, o globale, su $A$ se$$\forall\ x \in A,\ f(z) \geq f(x)$$.
\end{definition}

La definizione di minimo globale è analoga (basta sostituire il $\geq$ con $\leq$). 

Diamo adesso la definizione di massimo e minimo locale:

\vspace{0.5cm}
\begin{definition}[\textcolor{blue}{massimo locale}]
Sia $A \subseteq \mathbb{R}^n$ e sia $f: A \to \mathbb{R}$. Diremo che $z \in A$ è un punto di massimo locale su $A$ se
$$
\exists \delta > 0: \forall x \in A \cap B(z, \delta), f(z) \geq f(x)
$$
\end{definition}

\noindent analogamente si ottiene la definizione di minimo locale. \\ Data una funzione $f:A \to \mathbb{R}$ con $A \subset \mathbb{R}^n$, introdurremo la seguente notazione 
\begin{equation*}
\max_A{f} \qquad \min_A{f} 
\end{equation*}
per indicare, rispettivamente, il massimo e il minimo assoluto. 

Enunciamo il seguente teorema, valido in $\mathbb{R}^n$, che garantisce l'esistenza del massimo e del minimo di una funzione continua quando \emph{mappa} un insieme compatto ad un altro. 

\vspace{0.5cm}
\begin{theorem}[\textcolor{blue}{di Weierstrass}]
Sia $A \subset \mathbb{R}^n$ un insieme compatto e sia $f: A \to \mathbb{R}$ una funzione continua su $A$. Allora $\exists \max\limits_A{f}, \min\limits_A{f}$. 
\label{thm:weierstrass}
\end{theorem}

\noindent \textbf{Dimostrazione. }
Per caratterizzazione del $\sup{f(A)}$ sappiamo che esiste una successione $y_k$ tale che $y_k \in f(A) \, \forall k \geq 1$ e $y_k \stackrel{k \to +\infty}{\to} \sup f(A)$, ovvero esiste una successione a valori in $f(A)$ che tende all'estremo superiore di tale insieme. Siccome $\forall \, k \geq 1, \, y_k \in f(A)$ allora $\exists x_k \in A: f(x_k) = y_k \, \forall k \geq 1$. Ma allora noi sappiamo che, per al compattezza di $A$, che esiste una sottosuccessione $x_{\alpha (k)}$ tale che $x_{\alpha(k)} \to z \in A$, da cui si deduce che
$$
f(z) = \lim_{k \to +\infty} f(x_{\alpha(k)}) = \lim_{k \to +\infty} y_{\alpha(k)} = \sup_A{f}
$$
dove la prima uguaglianza segue dalla continuità di $f$ e la terza uguaglianza segue dal fatto che $y_{\alpha(k)}$ è una sottosuccessione estratta di $y_k$, ma siccome $y_k$ converge al $\sup\limits_A{f}$ allora anche $y_{\alpha(k)}$ vi convergerà.
\hfill $\square$

\vspace{0.5cm}
Un'altra buona proprietà delle funzioni continue è quello di preservare la connessione per archi di un insieme
\begin{definition}[\textcolor{blue}{insieme connesso per archi}]
	Un insieme $C \subseteq \mathbb{R}^n$ è connesso per archi se $\forall x, y \in C \exists \gamma: [0, 1] \to C$ continua tale che $\gamma(0) = x$ e $\gamma(1) = y$.
\end{definition}

\begin{example}
	Una palla è un insieme connesso per archi
\end{example}

\vspace{0.5cm}
\begin{definition}[\textcolor{blue}{insiemi convessi}]
	$C \subseteq \mathbb{R}^n$ è convesso $\iff \, \, \forall x, y \in C$ si ha che $[x, y] = \{ x + t(y-x) : t \in [0, 1] \} \subset C$.
\end{definition}

\begin{remark}
	In maniera spicciola, un insieme è convesso se, presi due punti, il segmento che li congiunge è costituito da punti tutti appartenenti all'insieme $C$. Le palle aperte e chiuse rimangono un esempio di insiemi convessi.
\end{remark}

\begin{remark}
	Dalla continuità del segmento segue banalmente che se $E \text{ convesso } \implies E \text{ connesso per archi}$. Il viceversa naturalmente non è valido: si potrebbe pensare all'insieme 
	$$
	A = \{(x, y) \in \mathbb{R}^2 : 1 \leq \sqrt{x^2 + y^2} \leq 2 \} = \mathbb{B}(0, 2) \setminus B(0,1)
	$$
	che corrisponde ad una corona circolare. E' connesso per archi siccome presi due punti possiamo trovare una curva che li "collega" tuttavia non è detto che, presi i due punti, sia possibile individuare un segmento che non passi nella zona che non fa parte di $A$.
\end{remark}

\noindent Dimostriamo adesso che la continuità preserva la connessione per archi di un insieme:

\begin{theorem}[C4]
	Sia $f: C \to \mathbb{R}^m$ continua e sia $C$ connesso per archi. Allora $f(C) \subset R^m$ è connesso per archi
\end{theorem}

\noindent \textbf{Dimostrazione. }
Siano $y, z \in f(C) \implies \exists x, u \in C$ tali che $f(x) = y$ e $f(u) = z$, pertanto $\exists \gamma : [0,1] \to C$ grazie alla connettività per archi di $C$ per cui $\gamma(0) = x$ e $\gamma(1) = u$. Ma allora la curva $f \circ \gamma: [0,1] \to f(C)$ è continua (siccome composizione di funzioni continue) e $(f \circ \gamma)(0) = y$ e $(f \circ \gamma)(1) = z$, dunque $f(C)$ è connesso per archi.
\hfill $\square$

\vspace{0.5cm}
\begin{cor}
$f: C \to \mathbb{R}$ continua, $C$ è connesso per archi $\implies f(C) \subset \mathbb{R}$ è un intervallo. 
\end{cor}

\noindent \textbf{Dimostrazione. }
Dal precedente teorema segue, naturalmente, che $f(C)$ è connesso per archi. Mostriamo che è un intervallo: dati $t, s \in f(C), \exists \Gamma: [0,1] \to f(C) \subset \mathbb{R}$ continua tale che $\Gamma(0) = t$ e $\Gamma(1) = s$. Dunque $[t,s] \subseteq \Gamma([0, 1]) \subseteq f(C) \implies f(C)$ è un intervallo.
\hfill $\square$

\vspace{0.5cm}
\noindent Da questo corollario è anche possibile anche dimostrare la seguente proposizione
\begin{prop}
Sia $A$ connesso per archi e sia $f: A \to \mathbb{R}$ continua. Allora

\begin{enumerate}[label=\protect\circled{\arabic*}]
	\item Se $\exists \max\limits_{A} f, \min\limits_{A} f$ allora $f(A) = [\min\limits_{A} f, \max\limits_{A} f]$
	\item Se $\not\exists \max\limits_{A} f, \exists \min\limits_{A} f$ allora $f(A)=[\min\limits_{A} f, \sup\limits_{A} f)$
\end{enumerate}
\end{prop}

\vspace{0.5cm}
\begin{exercise}
Mostrare la proposizione precedente
\end{exercise}

\noindent \textbf{Dimostrazione. }
Per il corollario sappiamo che $f(A)$ è un intervallo. \\
Se il minimo e il massimo della funzione sono ben definiti, allora l'immagine è banalmente contenuta fra il minimo e il massimo. \\
Per quanto riguarda il secondo caso, possiamo supporre, senza perdere di generalità, che $\not\exists \max\limits_{A} f$. Naturalmente, per caratterizzazione del $\sup$ abbiamo che $\forall \varepsilon > 0, \sup\limits_{A}{f} - \varepsilon \in f(A)$, dunque $\exists x \in A: f(x) = \sup\limits_{A}{f} - \varepsilon$, dunque $[\min\limits_{A} f, f(x)] \subseteq [\min\limits_A{f}, \sup\limits_{A}{f})$, dunque la tesi.
\hfill $\square$

\vspace{0.5cm}
\section{Accenni ai limiti in più variabili}
Per poter parlare dei limiti di funzioni a più variabili bisogna rivedere il concetto di punto di accumulazione, la cui definizione dev'essere leggermente modificata
\vspace{0.5cm}
\begin{definition}[\textcolor{blue}{punto di accumulazione}]
	Sia $A \subseteq \mathbb{R}^n e x_0 \in \mathbb{R}^n$, diremo che $x_0$ è un punto di accumulazione di $A$ se
	$$
		\forall r > 0, B(x, r) \cap A \setminus \{ x \} \neq \emptyset.
	$$
	Indichiamo con $D(A)$ l'insieme di tutti i punti di accumulazione di $A$.
\end{definition}
\vspace{0.5cm}
Possiamo adesso introdurre la definizione di limite:

\begin{definition}[\textcolor{blue}{definizione di limite}]
Sia data $f: A \to \mathbb{R}^k, A \subseteq \mathbb{R}^n$ e sia $x_0 \in D(A)$. Allora diciamo che $f$ tende a $v \in \mathbb{R}^k$ per $x$ che tende a $x_0$ e scriviamo che
$$
\lim_{x \to x_0} f(x) = v \iff \forall \varepsilon > 0 \exists \delta > 0 : \forall x \in B(x_0, \delta) \cap A \setminus \{ x_0 \}, |f(x) - v| < \varepsilon
$$
\end{definition}

\vspace{0.5cm}
\noindent Diamo adesso la definizione di limite di $f(x)$ che tende all'infinito:

\begin{definition}[\textcolor{blue}{limite che tende all'infinito}]
Sia data $f: A \to \mathbb{R}^k, A \subseteq \mathbb{R}^n$ illimitato. Diremo che $f$ tende all'infinito per $x$ che tende all'infinito e scriveremo che
$$
\lim_{|x| \to \infty} f(x) = \infty \iff \forall M > 0, \exists N > 0: \forall x \in A, |x| > N, |f(x)| > M
$$
\end{definition}

\begin{remark}
In generale tutte queste definizioni che stiamo dando funzionano in qualunque spazio metrico $(X,d)$ sostituendo al posto dell'usuale distanza euclidea definita in $\mathbb{R}^n$ la distanza $d$.
\end{remark}

\noindent Potremo scrivere anche le definizioni per gli altri due casi che ci restano, ma è banale e lo lasciamo come semplice esercizio teorico al lettore. \\

\begin{theorem}[\textcolor{blue}{caratterizzazione del limite per successioni}]
Sia $f: A \to \mathbb{R}^m, A \subseteq \mathbb{R}^n, q \in D(A) A, v \in \mathbb{R}^m$. Abbiamo che
$$
\lim_{x \to q} f(x) = v \iff (\forall \{ p_k \} \subseteq A \setminus \{ q \}, p_k \to q) \implies \lim_{k \to +\infty} f(p_k) = v
$$
\end{theorem}

\noindent \textbf{Dimostrazione. }
$\boxed{\Rightarrow}$: si osserva che, presa una qualunque successione $p_k$ che soddisfa le ipotesi della proposizione sulla destra, avremo che il limite $\lim\limits_{k \to +\infty} f(p_k)$ è un limite dove avviene la composizione della successione $p_k$ con la funzione $f$ e sono verificate tutte le ipotesi riguardo al cambio di variabile nel limite, dunque $\lim\limits_{k \to +\infty} f(p_k) = \lim\limits_{x \to q} f(x) = v$. \\
$\boxed{\Leftarrow}$: osserviamo che se procedessimo per assurdo, negando la definizione di limite, avremo che $\exists \varepsilon: \forall \delta > 0, \exists x \in B(q, \delta) \cap A \setminus \{ q \}: |f(x) - v| \geq \varepsilon$. Ma allora, restringendosi a degli intorni di $q$ via via sempre più piccoli, come $V_k = (q-\frac{1}{k}, q + \frac{1}{k})$ esisterebbe $x_k \in A \cap V_k : |f(x_k) - v| > \varepsilon$. Ma questo è un assurdo, siccome $x_k \to q$ e ma $f(x_k) \not\to v$ contraddicendo la nostra ipotesi iniziale.
\hfill $\square$

\vspace{0.5cm}
\noindent Da questo teorema importantissimo segue questo semplice corollario (la cui dimostrazione è omessa siccome segue direttamente dalla definizione di funzione continua)
\begin{cor}
	Sia $\Omega \subseteq \mathbb{R}^n$ aperto e $f: \Omega \to \mathbb{R}^m$. Allora valgono le seguenti proprietà:
	\begin{enumerate}[label=\protect\circled{\arabic*}]
		\item $f$ è continua in $x_0 \in \Omega \iff \lim\limits_{x \to q} f(x) = f(q)$;
		\item $f$ è continua in $\Omega \iff \, \, \forall q \in \Omega, \lim\limits_{x \to q} f(x) = f(q)$
	\end{enumerate}
\end{cor}

\vspace{0.5cm}
\begin{theorem}[\textcolor{blue}{del confronto}]
Siano $h, g, f: A \to \mathbb{R}, x_0 \in D(A) A, l \in \mathbb{R}$ e supponiamo che $\exists \delta > 0$ con
\[\begin{aligned}
h(x) \leq f(x) \leq g(x) \, \, \forall x \in A \cap B(x_0, \delta) \setminus \{ x_0 \}
\end{aligned}\]

Se 
$$\lim\limits_{x \to x_0} h(x) = \lim\limits_{x \to x_0} g(x) = l \implies \lim_{x \to x_0} f(x) = l$$
\end{theorem}

\noindent \textbf{Dimostrazione. }
La dimostrazione è analoga a quella di Analisi 1 con qualche accorgimento sugli insiemi di definizione
\hfill $\square$

\noindent Prima di procedere sulla parte del calcolo differenziale, diamo una breve introduzione ai simboli di Landau (i cosiddetti o-piccoli e O-grandi) che semplificano notevolmente il calcolo in più variabili:

\vspace{0.5cm}
\begin{definition}[\textcolor{blue}{o-piccolo}]
Siano $x_0 \in D(A) A, A \subseteq \mathbb{R}^n$ e $f: A \to \mathbb{R}^m$ e $g: A \to \mathbb{R}^k$. Diremo che $f$ è un o-piccolo di $g$ per $x$ che tende a $x_0$, in simboli
$$
f = o(g) \text{ per } x \to x_0
$$
se $$\forall \varepsilon > 0 \, \exists \delta > 0 : |f(x)| \leq \varepsilon |g(x)| \, \forall x \in A \cap B(x_0, \delta) \setminus \{ x_0 \}$$
\end{definition}

\vspace{0.5cm}
\begin{definition}[\textcolor{blue}{O-grande}]
Siano $x_0 \in D(A) A, A \subseteq \mathbb{R}^n$ e $f: A \to \mathbb{R}^m$ e $g: A \to \mathbb{R}^k$. Diremo che $f$ è un O-grande di $g$ per $x$ che tende a $x_0$, in simboli
$$
f = O(g) \text{ per } x \to x_0
$$
se $$\exists M > 0, \delta > 0 : |f(x)| \leq M|g(x)| \, \forall x \in A \cap B(x_0, \delta) \setminus \{ x_0 \} $$
\end{definition}

\vspace{0.5cm}
\begin{exercise}
Se $f$ è un o-piccolo di $g$ allora $f$ è O-grande di $g$
\end{exercise}

\noindent \textbf{Dimostrazione. }
Segue direttamente dalla definizione, siccome se $f = o(g)$ allora $\forall \varepsilon > 0 \exists \delta > 0: |f(x)| \leq \varepsilon |g(x)| \forall x \in A \cap B(x_0, \delta) \setminus \{ x_0 \}$, dunque fissando $M = \varepsilon > 0$ allora sappiamo che esiste un $\delta(\varepsilon)$ per cui
$$
|f(x)| \leq M |g(x)| \forall x \in A \cap B(x_0, \delta) \setminus \{ x_0 \}
$$ 
dunque risulta che $f$ è un O-grande di $g$, ovvero la tesi.
\hfill $\square$

\vspace{0.5cm}
\begin{prop}
Siano $A \subseteq \mathbb{R}^n, x_0 \in  A$ e $f: A \to \mathbb{R}^m$ e $\alpha > 0$. Allora abbiamo
$$
f(x) = o(|x-x_0|^\alpha) \text{ per } \, x \to x_0 \iff \lim_{x \to x_0} \frac{f(x)}{|x-x_0|^\alpha} = 0
$$
\end{prop}

\noindent \textbf{Dimostrazione. }
la dimostrazione è banale, siccome 

\[\begin{aligned}
&\forall \varepsilon > 0 \exists \delta > 0: |f(x)| \leq \varepsilon |x-x_0|^\alpha \, \forall x \in A \cap B(x_0, \delta) \setminus \{ x_0 \} \iff \forall \varepsilon > 0 \exists \delta > 0: \\
&\frac{|f(x)|}{|x-x_0|^\alpha} \leq \varepsilon \, \forall x \in A \cap B(x_0, \delta) \setminus \{ x_0 \} \iff \lim_{x \to x_0} \frac{f(x)}{|x-x_0|^\alpha} = 0
\end{aligned}\]
\hfill $\square$

\vspace{0.5cm}
\begin{prop}
Siano $A \subseteq \mathbb{R}^n, x_0 \in A, f: A \to \mathbb{R}$ e $\alpha > 0$. Allora abbiamo che $$f(x) = O(|x-x_0|^\alpha) \text{ per } x \to x_0 \iff \exists M, \delta > 0: \frac{f(x)}{|x-x_0|^\alpha} \leq M \, \forall x \in A \cap B(x_0, \delta) \setminus \{ x_0 \}$$
\end{prop}

\noindent \textbf{Dimostrazione. }
Analoga alla dimostrazione fatta sopra
\hfill $\square$

\noindent L'ultima proposizione si può anche riscrivere dicendo che se $f=O(|x-x_0|^\alpha)$ allora abbiamo necessariamente, per $x \to x_0$, che $\sup\limits_{A \cap B(x_0, \delta) \setminus \{ x_0 \}} \Bigg| \frac{f(x)}{|x-x_0|^{\alpha}} \Bigg| < +\infty$. \\
Osserviamo che
$$
\lim_{x \to x_0} f(x) = l \iff f(x) = l + o(1) \text{ per } x \to x_0
$$

\subsection{Proprietà}

Siano \( f_{\rho}: [0, 2\pi] \to \mathbb{R} \) e \( g_{\rho}: [0, 2\pi] \to \mathbb{R} \) due famiglie di funzioni dipendenti da \(\rho > 0\).

\begin{itemize}
	\item Se \( g_{\rho} = o(\rho^k) \), allora \( g_{\rho} = O(\rho^k) \).
	\item Se \( g_{\rho} = O(\rho^k) \), allora \( g_{\rho} = o(\rho^n) \) per ogni \( n = 0, 1, \ldots, k-1 \).
	\item Se \( f_{\rho} = O(\rho^k) \) e \( g_{\rho} = O(\rho^n) \), allora \( f_{\rho} g_{\rho} = O(\rho^{k+n}) \).
	\item Se \( f_{\rho} = O(\rho^k) \) e \( g_{\rho} = o(\rho^n) \), allora \( f_{\rho} g_{\rho} = o(\rho^{k+n}) \).
	\item Se \( f_{\rho} = O(\rho^k) \) e \( g_{\rho} = O(\rho^k) \), allora \( f_{\rho} + g_{\rho} = O(\rho^k) \).
	\item Se \( f_{\rho} = o(\rho^k) \) e \( g_{\rho} = o(\rho^k) \), allora \( f_{\rho} + g_{\rho} = o(\rho^k) \).
\end{itemize}

\subsection{Composizione con funzioni di una variabile}

Sia \( F: \mathbb{R} \to \mathbb{R} \) una funzione di una variabile.

\begin{itemize}
	\item Se \( F(t) = O(t^n) \) per un qualche \( n \geq 1 \) e \( g_{\rho} = o(\rho^k) \) per un \( k \geq 0 \), allora
	\[
	F \circ g_{\rho} = o(\rho^{kn}).
	\]
	\item Se \( F(t) = o(t^n) \) per un \( n \geq 0 \) e \( g_{\rho} = O(\rho^k) \) con \( k \geq 1 \), allora
	\[
	F \circ g_{\rho} = o(\rho^{kn}).
	\]
	\item Se \( F(t) = O(t^n) \) per un \( n \geq 0 \) e \( g_{\rho} = O(\rho^k) \) con \( k \geq 1 \), allora
	\[
	F \circ g_{\rho} = O(\rho^{kn}).
	\]
\end{itemize}

\subsection{Esempi}

\begin{itemize}
	\item \( \sin(\rho \sin \theta) = O(\rho) \)
	\item \( \sin(\rho \cos \theta) = O(\rho) \)
	\item \( \sin(\rho \sin \theta) - \rho \sin \theta = O(\rho^3) \)
	\item \( \cos(\rho \sin \theta) - 1 = O(\rho^2) \)
	\item \( e^{\rho \cos \theta} - 1 = O(\rho) \)
	\item \( e^{\rho \cos \theta} - 1 - \rho \cos \theta - \frac{1}{2} \rho \cos^2 \theta = O(\rho^3) \)
\end{itemize}